\documentclass[11pt]{article}
\usepackage[utf8]{inputenc}
\usepackage[margin=1in]{geometry}
\usepackage{amsmath, amssymb, amsthm}
\usepackage{graphicx}
\usepackage{hyperref}
\usepackage{enumitem}

% --- Custom Header Style from Scribe Reference ---
\newcommand{\scribe}[4]{
   \noindent
   \begin{center}
   \framebox{
      \vbox{
    \hbox to 6.28in { {\bf CSE400: Fundamentals of Probability in Computing} \hfill #2 }
       \vspace{4mm}
       \hbox to 6.28in { {\Large \hfill #1  \hfill} }
       \vspace{2mm}
       \hbox to 6.28in { {\it Lecturer: #3 \hfill Scribe: #4} }
      }
   }
   \end{center}
   \vspace*{4mm}
}

\begin{document}

\scribe{Inequalities and Tail Bounds}{March 24, 2026}{Dhaval Patel, PhD}{Prince}

\section{Introduction to Tail Bounds}
While the mean provides the average behavior of a random variable, tail bounds allow us to quantify the probability of "rare extreme events"[cite: 1391]. These are defined by the tail of a random variable $X$: $P\{X > x\}$[cite: 1373]. Tail bounds are essential in computing for analyzing scenarios like jobs delayed over 24 hours or packets dropped due to full buffers[cite: 1375].

\section{Chebyshev's Inequality}
Chebyshev's Inequality is more powerful than Markov's Inequality because it utilizes both the \textbf{mean} and the \textbf{variance} to control the deviation from the mean[cite: 1806, 1822].

\subsection{Definition}
For any random variable $X$ with mean $\mu$ and variance $\sigma^2$, and for any $a > 0$:
\begin{equation}
P(|X - \mu| \ge a) \le \frac{Var(X)}{a^2}
\end{equation}
[cite: 1825, 1846].

\subsection{Proof of Chebyshev's Inequality}
The proof relies on applying Markov's Inequality to the non-negative random variable $(X - \mu)^2$[cite: 1847]:
\begin{enumerate}
    \item Start with the tail probability: $P(|X - \mu| \ge a)$.
    \item Square both sides of the inequality: $P((X - \mu)^2 \ge a^2)$[cite: 1848].
    \item Apply Markov's Inequality ($P(Y \ge k) \le \frac{E[Y]}{k}$):
    \begin{equation}
    P((X - \mu)^2 \ge a^2) \le \frac{E[(X - \mu)^2]}{a^2}
    \end{equation}
    \item By definition, $E[(X - \mu)^2] = Var(X)$, thus:
    \begin{equation}
    P(|X - \mu| \ge a) \le \frac{Var(X)}{a^2}
    \end{equation}
\end{enumerate}
[cite: 1849].

\subsection{Solved Example: Exam Grades}
\textbf{Problem:} Average grade $\mu = 70$, standard deviation $\sigma = 5$. Find the bound for the fraction of students receiving an "A" (cutoff 90)[cite: 1930].
\begin{itemize}
    \item \textbf{Step 1:} Identify $\mu = 70$ and $\sigma^2 = 25$[cite: 1934].
    \item \textbf{Step 2:} $X \ge 90 \implies |X - 70| \ge 20$[cite: 1945].
    \item \textbf{Step 3:} $P(|X - 70| \ge 20) \le \frac{25}{20^2} = \frac{25}{400} = \frac{1}{16}$[cite: 1972, 1988].
    \item \textbf{Result:} At most 6.25\% of students receive an "A", which is much tighter than the Markov bound of 77.8\%[cite: 1985, 1988].
\end{itemize}

\section{Chernoff Bound}
The Chernoff bound provides the tightest tail bounds by using the \textbf{Moment Generating Function} ($e^{tX}$)[cite: 2071, 2122].

\subsection{Core Trick and Form}
The transformation magnifies large values of $X$ exponentially to create a tighter bound[cite: 2079, 2108].
\begin{equation}
P(X \ge a) = P(e^{tX} \ge e^{ta}) \le \frac{E[e^{tX}]}{e^{ta}} \text{ for } t > 0
\end{equation}
[cite: 2162, 2170]. The final form is minimized over $t$:
\begin{equation}
P(X \ge a) \le \min_{t>0} \left\{ \frac{E[e^{tX}]}{e^{ta}} \right\}
\end{equation}
[cite: 2186].

\subsection{Poisson Tail Bound Derivation}
For $X \sim Poisson(\lambda)$, the MGF is $E[e^{tX}] = e^{\lambda(e^t - 1)}$[cite: 2246].
Setting $t = \ln(a/\lambda)$ to minimize the exponent[cite: 2259]:
\begin{equation}
P(X \ge a) \le e^{a - \lambda} \left( \frac{\lambda}{a} \right)^a
\end{equation}
[cite: 2270].

\section{Jensen's Inequality}
Jensen's Inequality relates the value of a convex function of an integral to the integral of the convex function[cite: 2544].

\subsection{Convex Functions}
A function $g(x)$ is convex if $g(\alpha x_1 + (1-\alpha)x_2) \le \alpha g(x_1) + (1-\alpha)g(x_2)$[cite: 2440]. Mathematically, $g''(x) \ge 0$[cite: 2461].

\subsection{Theorem}
If $g(x)$ is convex on an interval $S$ and $X$ takes values in $S$, then:
\begin{equation}
g(E[X]) \le E[g(X)]
\end{equation}
[cite: 2546]. Interpretation: Averaging first gives a smaller value than applying the function first[cite: 2555, 2556].

\section{Summary for Review}
\begin{table}[h]
\centering
\begin{tabular}{|l|l|l|l|}
\hline
\textbf{Inequality} & \textbf{Information Used} & \textbf{Bound Type} & \textbf{Tightness} \\ \hline
Markov              & Mean                      & Polynomial ($1/a$)  & Loose [cite: 1622] \\ \hline
Chebyshev           & Mean + Variance           & Polynomial ($1/a^2$)& Moderate [cite: 1872] \\ \hline
Chernoff            & Full Distribution (MGF)   & Exponential ($e^{-a}$)& Strongest [cite: 2204] \\ \hline
\end{tabular}
\end{table}
\textbf{Golden Rule:} More statistical information leads to tighter bounds, which results in lower hardware costs in system design.

\end{document}
